\documentclass{article}
\usepackage[a4paper, left=1.5cm, right=1.5cm, top=1cm, bottom=2cm]{geometry}

\usepackage{boldline}
\usepackage{tikz,tcolorbox}
\usepackage{amsmath}
\usepackage[table,xcdraw]{xcolor}
\usepackage{listings}
\usepackage{array,multirow} % For customizing tables
\usepackage{booktabs} % For better horizontal lines
\usepackage{makecell}
\setlength{\parindent}{0pt}
\usepackage{siunitx}
\usepackage{tkz-tab}
\usepackage{amssymb}
\usepackage{amsmath}
\usepackage{enumitem}

\usepackage{caption}
\usepackage{float}


\newcommand{\exer}[1]{
  \section*{Exercice #1}
  \vspace{-0.5cm}
  \noindent\rule{\textwidth}{0.5pt}%
}

\newcommand{\tit}[1]{
\begin{center}
    \Large{\textbf{{#1}}}
\end{center}
}

\definecolor{commentgray}{HTML}{676160}
\definecolor{messagegreen}{HTML}{17B867}
\definecolor{myblue}{HTML}{10C2C4}

\tcbuselibrary{skins, breakable, theorems}
\usepackage{booktabs} 

\newtcolorbox{prettyBox}[2]{
  enhanced,
  colback=white!90!#2,   % Background color based on the second parameter (color)
  colframe=#2!60!black,  % Frame color based on the second parameter (color)
  coltitle=white,        % Title color (white)
  fonttitle=\bfseries\Large,
  title=#1,              % Title from the first parameter
  boxrule=1mm,
  arc=0.5mm,
  drop shadow=#2!35!gray, % Drop shadow color based on the second parameter (color)
}



\lstdefinestyle{cmd}{
 basicstyle=\ttfamily,
 backgroundcolor=\color{lightgray!20},
 frame=single
}

\lstdefinestyle{pythonstyle}{
    language=python,                    % Language set to Python
    basicstyle=\ttfamily\footnotesize,   % Change basic font size
    keywordstyle=\color{blue}\bfseries, % Different keyword style
    stringstyle=\color{red},         % Different string color
    commentstyle=\color{green!60!black}\itshape, % Adjust comment color
    numbers=left,                       % Line numbers on the left
    numberstyle=\tiny\color{gray},      % Smaller number font and color
    stepnumber=1,                       % Number each line
    frame=single,                       % Single frame around code
    tabsize=4,                          % Adjust tab size
    showstringspaces=false,             % Do not show spaces in strings
    captionpos=b,% Position of caption
    breaklines=true,
    inputencoding=utf8
}




\begin{document}
\tit{TD N\(^{\boldsymbol{\circ}}\)\hspace{0.1cm}2}
\exer{1}

\vspace{0.25cm}


\begin{prettyBox}{K-Means}{myblue}
    K-means est un algorithme non supervisé de clustering. Il prend comme paramètre le nombre de 
    clusters $K$ et $K$ centroides, et à la fin il nous donne les centroides finaux avec leurs points 
    formant les clusters.
\end{prettyBox}

\vspace{0.15cm}

\begin{prettyBox}{Algo}{myblue}
    \begin{enumerate}
        \setlength{\itemsep}{0pt}
        \item Choisir le nombre de clusters $K$.
        \item Choisir $K$ centroids.
        \item Calculer la distance entre chaque point et les centroids.
        \item Affecter chaque point au centroide le plus proche.
        \item Calculer les nouveaux centroides de chaque cluster en prenant la moyenne des points.
        \item Si les nouveaux clusters sont différents des anciens, revenir à l'étape 3 sinon, l'algorithme est terminé.
    \end{enumerate}
\end{prettyBox}

\vspace{0.15cm}
\begin{prettyBox}{Distance Euclidienne}{myblue}
    \[ \sqrt{\sum_{i = 1}^M x_i^2} \]
    \begin{itemize}
        \setlength{\itemsep}{0pt}
        \item $M$ : nombre d'attributs
        \item $x_i$ : ième attribut
    \end{itemize}
\end{prettyBox}

\vspace{0.15cm}
\begin{prettyBox}{SSE(Somme des Carrés des Erreurs)}{myblue}

   SSE est utilisée pour quantifier la cohésion interne (le degré d'appartenance des points à leur cluster). Elle est définie comme la somme des carrés des distances entre le
   point moyen (centroide) et chaque point du cluster.

    \[\sum_{i=1}^K \sum_{P \in S_i } D(P,C_i)^2\]

    \begin{itemize}
        \setlength{\itemsep}{0pt}
        \item $K$ : nombre de clusters.
        \item $P$ : instance (Point).
        \item $S_i$ : ième cluster.
        \item $C_i$ : ième centroide.
        \item $D$ : La distance.
    \end{itemize}
\end{prettyBox}

\newpage

On a comme data-set :

\begin{table}[h]
\centering

\renewcommand{\arraystretch}{1.5} 
\setlength{\tabcolsep}{12pt}      

\begin{tabular}{|l|c|c|} 
\hline
\textbf{Client} & \textbf{Dépenses (X) (en milliers d'€)} & \textbf{Fréquence (Y) (sessions/an)} \\ \hline
A        & 2          & 8        \\ \hline
B          & 3           & 7             \\ \hline
C      & 4      & 5           \\ \hline
D        & 8          & 2        \\ \hline
E          & 7           & 3             \\ \hline
F      & 6     & 4          \\ \hline

\end{tabular}
\end{table}

On veut segmenter les clients en 2 groupe ($K=2$) avec les centroides suivant :\\[0.1cm]
$C_1 = A = (2,8) \quad C_2 = D = (8,2) $


\vspace{0.75cm}
Calculer la distance Euclidienne de chaque point aux deux centres initiaux ($C_1 , C_2$) :

\begin{table}[h]
\centering

\renewcommand{\arraystretch}{2.25} 
\setlength{\tabcolsep}{3pt}      

\begin{tabular}{|c|c|c|c|c|} 
\hline
\textbf{Client($P$)} & \textbf{Coordonnées} & $\mathbf{D(P,C_1)}$ & $\mathbf{D(P,C_2)}$ & Cluster Attribue  \\ \hline
A        & (2,8)          & $\sqrt{(2-2)^2 + (8-8)^2} = 0$  & $\sqrt{(2-8)^2 + (8-2)^2} = \sqrt{72} = 8.49$    & $C_1$  \\ \hline
B          & (3,7) & $\sqrt{(3-2)^2 + (7-8)^2} = \sqrt{2} = 1.41$  & $\sqrt{(3-8)^2 + (7-2)^2} = \sqrt{50} = 7.07$    & $C_1$  \\ \hline
C & (4,5) & $\sqrt{(4-2)^2 + (5-8)^2} = \sqrt{13} = 3.61$  & $\sqrt{(4-8)^2 + (5-2)^2} = \sqrt{25} = 5$    & $C_1$  \\ \hline
D        & (8,2)& $\sqrt{(8-2)^2 + (2-8)^2} = \sqrt{72} = 8.49$  & $\sqrt{(8-8)^2 + (2-2)^2} = 0$    & $C_2$  \\ \hline     
E          & (7,3)& $\sqrt{(7-2)^2 + (3-8)^2} = \sqrt{50} = 7.07$  & $\sqrt{(7-8)^2 + (3-2)^2} = \sqrt{2} = 1.41$    & $C_2$  \\ \hline           
F      & (6,4) & $\sqrt{(6-2)^2 + (4-8)^2} = \sqrt{32} = 5.66$  & $\sqrt{(6-8)^2 + (4-2)^2} = \sqrt{8} = 2.83$    & $C_2$  \\ \hline           

\end{tabular}
\end{table}

\vspace{0.75cm}

Attribution des points au clusters :
\begin{itemize}
    \setlength{\itemsep}{0pt}
    \item \textbf{Cluster 1} ($S_1$) : $A , B , C$
    \item \textbf{Cluster 2} ($S_2$) : $D , E , F $
\end{itemize}

\vspace{0.75cm}

Calcule des nouveaux centroides :

\begin{align*}
&C^{'}_1 = \left(\dfrac{x_A+x_B+x_C}{3} , \dfrac{y_A+y_B+y_C}{3} \right) = \left(\dfrac{2+3+4}{3} , \dfrac{8+7+5}{3} \right)  = (3,6.67)\\[0.15cm]
&C^{'}_2 = \left(\dfrac{x_D+x_E+x_F}{3} , \dfrac{y_D+y_E+y_F}{3} \right) = \left(\dfrac{8+7+6}{3} , \dfrac{2+3+4}{3} \right)  = (7,3)
\end{align*}

\vspace{0.45cm}

$C^{'}_1 = (3,6.67)$ est le nouveau centroide du cluster $S_1$ des points $ A , B , C$ \\[0.15cm] 
$C^{'}_2 = (7,3)$ est le nouveau centroide du cluster $S_2$ des points $ D , E , F$ 

\newpage 

Évaluer la qualité de la première itération en calculant la Somme des Carrés des Erreurs (SSE) : 

\[\sum_{i=1}^{K=2} \sum_{P \in S_i } D(P,C^{'}_i)^2 = \sum_{P \in S_1 } D(P,C^{'}_1)^2 + \sum_{P \in S_2 } D(P,C^{'}_2)^2  \]


\begin{table}[h]
\centering


\renewcommand{\arraystretch}{2.25} 
\setlength{\tabcolsep}{3pt}      

\begin{tabular}{|c|c|c|c|} 
\hline
\textbf{Client($P$)} & \textbf{Coordonnées} & Cluster Attribue  & $(\mathbf{D(P,C^{'}_i)})^2$  \\ \hline
A        & (2,8)      & $C^{'}_1 = (3,6.67)$    &  $(2-3)^2 + (8-6.67)^2  = 2.77$     \\ \hline
B          & (3,7) & $C^{'}_1 = (3,6.67)$ &   $(3-3)^2 + (7-6.67)^2 = 0.11$     \\ \hline
C & (4,5)  & $C^{'}_1 = (3,6.67)$  & $ (4-3)^2 + (5-6.67)^2 = 3.79$      \\ \hline
D        & (8,2) & $C^{'}_2 = (7,3)$ &   $(8-7)^2 + (2-3)^2 = 2$    \\ \hline     
E          & (7,3) & $C^{'}_2 = (7,3)$  & $(7-7)^2 + (3-3)^2 = 0$ \\ \hline           
F      & (6,4) & $C^{'}_2 = (7,3)$  & $(6-7)^2 + (4-3)^2 = 2$   \\ \hline           

\end{tabular}
\end{table}

\begin{align*}
    \textbf{\text{SSE}} & = \sum_{P \in S_1 } D(P,C^{'}_1)^2 + \sum_{P \in S_2 } D(P,C^{'}_2)^2\\[0.1cm]
                        & = (2.77 + 0.11 + 3.79) + (2+0+2)\\[0.1cm]
                        & =  6.67 + 4\\[0.1cm]
                        & = 10.67
\end{align*}


\begin{table}[h]
\centering

\renewcommand{\arraystretch}{2.25} 
\setlength{\tabcolsep}{0.75pt}      

\begin{tabular}{|c|c|c|c|c|} 
\hline
\textbf{Client($P$)} & \textbf{Coordonnées} & $\mathbf{D(P,C^{'}_1)}$ & $\mathbf{D(P,C^{'}_2)}$ & Cluster Attribue  \\ \hline
A        & (2,8)          & $\sqrt{(2-3)^2 + (8-6.67)^2} = \sqrt{2.7689} =1.66 $  & $\sqrt{(2-7)^2 + (8-3)^2} = \sqrt{50} = 7.07$    & $C^{'}_1$  \\ \hline
B          & (3,7) & $\sqrt{(3-3)^2 + (7-6.67)^2} = \sqrt{0.1089} = 0.33$  & $\sqrt{(3-7)^2 + (7-3)^2} = \sqrt{32} = 5.66$    & $C^{'}_1$  \\ \hline
C & (4,5) & $\sqrt{(4-3)^2 + (5-6.67)^2} = \sqrt{3.7889} = 1.95$  & $\sqrt{(4-7)^2 + (5-3)^2} = \sqrt{13} = 3.61$    & $C^{'}_1$  \\ \hline
D        & (8,2)& $\sqrt{(8-3)^2 + (2-6.67)^2} = \sqrt{46.8089} = 8.81$  & $\sqrt{(8-7)^2 + (2-3)^2} = \sqrt{2} =1.41$    & $C^{'}_2$  \\ \hline     
E          & (7,3)& $\sqrt{(7-3)^2 + (3-6.67)^2} = \sqrt{29.4689} = 5.43$  & $\sqrt{(7-7)^2 + (3-3)^2} = 0$    & $C^{'}_2$  \\ \hline           
F      & (6,4) & $\sqrt{(6-3)^2 + (4-6.67)^2} = \sqrt{16.1289} = 4.02$  & $\sqrt{(6-7)^2 + (4-3)^2} = \sqrt{2} = 1.41$    & $C^{'}_2$  \\ \hline           

\end{tabular}
\end{table}


\vspace{0.45cm}

On remarque que les points ne changent pas de clusters donc l'algo converge avec les centroides $C^{'}_1 , C^{'}_2$ :

\begin{itemize}
    \setlength{\itemsep}{0pt}
    \item \textbf{Cluster 1 (Clients à forte fréquence, dépenses modérées)} : \{$A, B, C$\}
    \item \textbf{Cluster 2 (Clients à faible fréquence, grosses dépenses)} : \{$D, E, F$\}
\end{itemize}

%\begin{figure}[H] 
 % \centering 
  %\includegraphics[width=0.5\textwidth]{Q.png} 
%\end{figure}

\end{document}
